\section{Ground space of non-injective MPS}\label{sec:degenerate_ground_space}

For an MPS with more than one block in canonical form, after blocking to
block-injective canonical form, the physical--virtual correspondence is
restricted to block-diagonal boundary conditions.  In the notation of
\cite[lines~2113--2117]{Cirac2021Matrix}, set
\[
    \widetilde A^\omega:=(A_1^\omega,\ldots,A_g^\omega)
    \in\prod_{j=1}^g \MN{D_j}.
\]
Then block injectivity is the statement
\[
    \forall X=(X_j)_{j=1}^g\in\bigoplus_{j=1}^g \MN{D_j},
    \quad
    \exists (c_\omega(X))_\omega,
    \quad
    X=\sum_\omega c_\omega(X)\widetilde A^\omega.
\]
Equivalently, with the same coefficient family,
\[
    \forall j,\qquad
    X_j=\sum_\omega c_\omega(X)A_j^\omega .
\]
The parent-Hamiltonian theorem in \cite[lines~2121--2128]{Cirac2021Matrix} is
\[
    N\ge L_0+1
    \quad\Longrightarrow\quad
    \ker H_N=\spn\{V^{(N)}(A_j):j=1,\ldots,g\}.
\]
The statements below record the corresponding cyclic-chain ground-space
formulation; the passage to $\ker H_N$ is kept separate.

\begin{definition}[Chain ground space for the block-diagonal tensor]
    \label{def:parent_hamiltonian_ground_space_bnt}
    \notready
    \uses{def:chain_ground_space}
    For a basis of normal tensors $A=(A_j)_{j=1}^g$ and nonzero weights
    $\mu_j$, let
    \[
        B := \bigoplus_{j=1}^g \mu_j A_j .
    \]
    The periodic chain ground space considered in this section is
    \[
        \mathcal G_{N,L}(B).
    \]
\end{definition}

\begin{lemma}[Block-diagonal chain ground space]
    \label{lem:parent_hamiltonian_ground_space_eq}
    \notready
    \uses{def:parent_hamiltonian_ground_space_bnt}
    By definition, for $B=\bigoplus_{j=1}^g \mu_j A_j$ one has
    \[
        \mathcal G_{N,L}(B)
        =
        \mathcal G_{N,L}\!\left(\bigoplus_{j=1}^g \mu_j A_j\right).
    \]
\end{lemma}

\begin{definition}[Vector space generated by a basis of normal tensors]
    \label{def:bnt_span}
    \notready
    For a basis of normal tensors $A_1,\ldots,A_g$, the vector space generated
    by their length-$N$ matrix product vectors is
    \[
        \spn\{V^{(N)}(A_j): j=1,\ldots,g\}.
    \]
\end{definition}

\begin{definition}[Finite blockwise physical--virtual correspondence]
    \label{def:blockwise_product_word_span}
    \lean{MPSTensor.WordTupleSpanTop}
    \leanok
    \uses{rem:word_tuple_span_top}
    For a fixed blocked length $m$, the finite-length blockwise
    physical--virtual correspondence is
    \[
        \spn\{(A_1^\omega,\ldots,A_g^\omega): |\omega|=m\}
        =
        \prod_{j=1}^g \MN{D_j}.
    \]
    Equivalently,
    \[
        \forall X=(X_j)_j\in\bigoplus_{j=1}^g \MN{D_j},
        \quad
        \exists (c_\omega(X))_{|\omega|=m},
        \quad
        (X_j)_{j=1}^g
        =
        \sum_{|\omega|=m} c_\omega(X)(A_1^\omega,\ldots,A_g^\omega).
    \]
    This is the fixed-length form of the block-injective correspondence of
    \cite[lines~2113--2117]{Cirac2021Matrix}.
\end{definition}

\begin{lemma}[Each normal-tensor vector lies in the chain ground space]
    \label{lem:bnt_mpv_mem_ground_space}
    \notready
    \uses{def:parent_hamiltonian_ground_space_bnt, lem:mpv_mem_chain_ground_space}
    If $1 < L$ and $N \ge L + 1$, then for each normal tensor $A_j$ in the basis
    the vector $V^{(N)}(A_j)$ lies in
    $\mathcal G_{N,L}(\bigoplus_k \mu_k A_k)$.
\end{lemma}

\begin{proof}
    \notready
    Write $\Gamma_N^B$ for the length-$N$ boundary-to-state map of the tensor
    $B$, and let $\iota_j(Z)$ be the block-diagonal matrix with $Z$ in the
    $j$-th block and $0$ elsewhere. The direct-sum evaluation gives
    \[
        \Gamma_N^{\oplus_k\mu_k A_k}\!\left(\iota_j(\mu_j^{-N}X)\right)
        =
        \Gamma_N^{A_j}(X).
    \]
    Taking $X=I$ gives the vector $V^{(N)}(A_j)$, and
    Lemma~\ref{lem:mpv_mem_chain_ground_space} applied to the individual block
    gives its membership in $\mathcal G_{N,L}(\bigoplus_k\mu_k A_k)$.
\end{proof}

\begin{lemma}[Inclusion of a block ground space]
    \label{lem:chain_ground_space_block_le_assembled}
    \lean{MPSTensor.chainGroundSpace_block_le_toTensorFromBlocks}
    \leanok
    \uses{def:chain_ground_space, lem:local_ground_space_block_le_assembled}
    For each block index $j$ with $\mu_{j} \neq 0$, the chain ground space of the block
    $A_j$ embeds into the chain ground space of the direct-sum tensor:
    \[
        \mathcal G_{N,L}(A_{j})
        \subseteq
        \mathcal G_{N,L}\!\left(\bigoplus_k \mu_k A_k\right).
    \]
\end{lemma}

\begin{proof}
    \leanok
    The cyclic-window definition of the chain ground space reduces the inclusion to
    the local block-embedding statement that each ground vector of $A_j$ on the
    window of length $L$ lifts to the direct sum by the identity
    \[
        \Gamma_L^{\oplus_k\mu_k A_k}\!\left(\iota_j(\mu_j^{-L}X)\right)
        =
        \Gamma_L^{A_j}(X).
    \]
    Here $\iota_j$ inserts a matrix in the $j$-th diagonal block and has zero
    off that block. Thus $\iota_j(\mu_j^{-L}X)$ has exactly the local vector of
    the $j$-th block on that window.
\end{proof}

\begin{lemma}[Forward inclusion from the block ground spaces]
    \label{lem:isup_chain_ground_space_block_le_assembled}
    \lean{MPSTensor.iSup_chainGroundSpace_block_le_toTensorFromBlocks}
    \leanok
    \uses{def:chain_ground_space, lem:chain_ground_space_block_le_assembled}
    Assume $\mu_{j} \neq 0$ for every block index $j$. The linear sum of the
    blockwise periodic-chain ground spaces is contained in the direct-sum
    tensor's periodic-chain ground space:
    \[
        \sum_{j} \mathcal G_{N,L}(A_{j})
        \subseteq
        \mathcal G_{N,L}\!\left(\bigoplus_k \mu_k A_k\right),
    \]
    where $j$ ranges over all block indices.
\end{lemma}

\begin{proof}
    \leanok
    Take the linear sum of the inclusions
    \[
        \mathcal G_{N,L}(A_j)
        \subseteq
        \mathcal G_{N,L}\!\left(\bigoplus_k \mu_k A_k\right)
    \]
    over all block indices $j$.
\end{proof}

\begin{lemma}[Forward inclusion for the direct-sum tensor]
    \label{lem:isup_chain_ground_space_block_le_parent_ground_space}
    \notready
    \uses{def:parent_hamiltonian_ground_space_bnt,
        lem:isup_chain_ground_space_block_le_assembled}
    Let $A^i=\bigoplus_j\mu_j A_j^i$ be in block-injective canonical form,
    where $A_1,\ldots,A_g$ is a basis of normal tensors.  The linear sum of the
    blockwise periodic-chain ground spaces is contained in the periodic chain
    ground space of the block-diagonal tensor:
    \[
        \sum_{j} \mathcal G_{N,L}(A_{j})
        \subseteq
        \mathcal G_{N,L}\!\left(\bigoplus_k \mu_k A_k\right).
    \]
\end{lemma}

\begin{proof}
    \notready
    The block-injective canonical-form hypothesis supplies $\mu_{j} \neq 0$ for
    every $j$, and the block-diagonal chain ground space is
    $\mathcal G_{N,L}(\bigoplus_j \mu_j A_j)$ by definition.
\end{proof}
